% ============================================================
% APÊNDICE C — Parâmetros de controle e limitadores operacionais
% ============================================================

% \section{Parâmetros de controle e limitadores operacionais}
\label{ap_apendiceC_parametros_controle}

Este apêndice consolida os parâmetros numéricos utilizados nas leis de controle e nos limitadores operacionais do perseguidor e do \gls{mr}.

\section{Tempo de simulação e discretizações auxiliares}
\label{apC_tempo_dt}

\begin{table}[H]
\centering
\caption{Parâmetros da simulação.}
\label{ttab_apC_tempo_dt}
\begin{tabular}{lr}
\toprule
\textbf{Parâmetro} & \textbf{Valor} \\
\midrule
Tempo inicial $t_0$ [s]          & 0 \\
Duração total [min]              & 15 \\
Tempo final [s]                  & 900 \\
Passo de integração $dt$ [s]     & 0,01 \\
\bottomrule
\end{tabular}
\end{table}


%%%%%%%%%%%%%%%%%%

\section{Saturações e limites da dinâmica de Lagrange}
\label{apC_saturacoes_lagrange}

\begin{table}[H]
\centering
\caption{Limites da espaçonave e das juntas no modo operacional.}
\label{ttab_apC_saturacoes_lagrange}
\begin{tabular}{lr}
\toprule
\textbf{Limite} & \textbf{Valor} \\
\midrule
Saturação de aceleração da espaçonave $\,\ddot{(\cdot)}_{\text{perseg,max}}$ & 2 \\
Saturação de velocidade da espaçonave $\,\dot{(\cdot)}_{\text{perseg,max}}$  & 2 \\
Saturação de aceleração das juntas $\,\ddot{q}_{\max}$                       & 2 \\
Saturação de velocidade das juntas $\,\dot{q}_{\max}$                        & 2 \\
Saturação de posição das juntas $\,q_{\max}$ [rad]                           & 6,28 \\
\bottomrule
\end{tabular}
\end{table}

%%%%%%%%%%%%%%%%%%%%%%%%

\section{Controle PD de atitude do perseguidor (modo não operacional)}
\label{apC_pd_euler}

\begin{table}[H]
\centering
\caption{Parâmetros do controle de atitude no modo não operacional.}
\label{ttab_apC_pd_euler}
\begin{tabular}{lr}
\toprule
\textbf{Parâmetro} & \textbf{Valor} \\
\midrule
Ganho de torque escalar & 50 \\
$\vec{PD}_{\tau,\text{Euler}}$ (3 eixos) & $[50;50;50]^T$ \\
Ganho proporcional & 10 \\
Ganho derivativo & 5 \\
$\vec{K}_{p,\text{Euler}}$ (3 eixos) & $[10;10;10]^T$ \\
$\vec{K}_{d,\text{Euler}}$ (3 eixos) & $[5;5;5]^T$ \\
\bottomrule
\end{tabular}
\end{table}

%%%%%%%%%%%%%%%%%%%%%%%%%%%%%%

\section{Controle PD de atitude da espaçonave (modo operacional).}
\label{apC_pd_lagrange}


\begin{table}[H]
\centering
\caption{Parâmetros do controle PD de atitude no modo operacional.}
\label{ttab_apC_pd_lagrange}
\begin{tabular}{lr}
\toprule
\textbf{Parâmetro} & \textbf{Valor} \\
\midrule
Ganho de torque escalar & 50 \\
$\vec{PD}_{\tau,\text{Lagrange}}$ (3 eixos) & $[50;50;50]^T$ \\
Ganho proporcional & 30 \\
Ganho derivativo & 15 \\
$\vec{K}_{p,\text{Lagrange}}$ (3 eixos) & $[30;30;30]^T$ \\
$\vec{K}_{d,\text{Lagrange}}$ (3 eixos) & $[15;15;15]^T$ \\
\bottomrule
\end{tabular}
\end{table}

%%%%%%%%%%%%%%%%%%%%

\section{Controle do manipulador robótico}
\label{apC_controle_mr}

\begin{table}[H]
\centering
\caption{Parâmetros do controle articular do \gls{mr} (3R).}
\label{ttab_apC_controle_mr}
\begin{tabular}{p{0.70\linewidth}r}
\toprule
\textbf{Parâmetro} & \textbf{Valor} \\
\midrule
Ganho proporcional escalar (articular) $K_{p,\text{MR}}$ & 1 \\
Ganho derivativo escalar (articular) $K_{d,\text{MR}}$   & 1 \\
Ganho proporcional do controle articular $\vec{K}_{p,\text{MR}}$ & $[3;3;3]^T$ \\
Ganho derivativo do controle articular $\vec{K}_{d,\text{MR}}$   & $[1,5;1,5;1,5]^T$ \\
Saturação de torque articular $\vec{\tau}_{\max}$                & $[2;2;2]^T$ \\
Saturação escalar de torque (articular) $\tau_{\text{sat}}$      & 2 \\
Ganho escalar de torque (articular) $\tau_{\text{nom}}$          & 1 \\
Termo de \textit{feedforward} $\vec{K}_{\text{FF}}$              & $[0;0;0]^T$ \\
\bottomrule
\end{tabular}
\end{table}



%%%%%%%%%%%%%%%%%%%%%%%%%%

\begin{table}[H]
\centering
\caption{Parâmetros da cinemática inversa e do mapeamento geométrico do \gls{mr}.}
\label{ttab_apC_ik_mr}
\begin{tabular}{lr}
\toprule
\textbf{Parâmetro} & \textbf{Valor} \\
\midrule
Ganho de posição & 1 \\
Ganho de alinhamento & 0,5 \\
Limite de velocidade articular (IK) [rad/s] & 0,1 \\
Damping (IK) & 2 \\
\bottomrule
\end{tabular}
\end{table}



%%%%%%%%%%%%%%%%%%
% \begin{table}[H]
% \centering
% \caption{Parâmetros do controle por \textit{twist} (quando habilitado).}
% \label{ttab_apC_twist}
% \begin{tabular}{lr}
% \toprule
% \textbf{Parâmetro} & \textbf{Valor} \\
% \midrule
% Ganho angular & 3 \\
% Ganho lateral & 2 \\
% Ganho axial & 2 \\
% $\epsilon$ lateral & 0,1 \\
% $\epsilon$ angular [graus] & 0,2 \\
% Saturação de velocidade de referência & 1 \\
% \bottomrule
% \end{tabular}
% \end{table}


%%%%%%%%%%%%%%%%%%%%%%%
\section{Critério de ativação do modo operacional do \gls{mr}}
\label{apC_ativacao_modo_op}

\begin{table}[H]
\centering
\caption{Parâmetros da lógica de ativação do modo operacional do \gls{mr} na missão.}
\label{ttab_apC_modo_op}
\begin{tabular}{lr}
\toprule
\textbf{Parâmetro} & \textbf{Valor} \\
\midrule
Fator do critério radial & 1,2 \\
Tempo mínimo de permanência [s] & 15 \\
\bottomrule
\end{tabular}
\end{table}

%%%%%%%%%%%%%%%%%%%%%%%%%%%%%%%%%
\section{Resumo de modos (missão)}
\label{apD_resumo_modos}

\begin{table}[H]
\centering
\caption{Resumo dos modos principais da missão e subsistemas habilitados.}
\label{ttab_apD_modos_missao}
\small
\begin{tabular}{p{0.10\linewidth}p{0.30\linewidth}p{0.56\linewidth}}
\toprule
\textbf{Modo} & \textbf{Descrição} & \textbf{Subsistemas habilitados} \\
\midrule
0 & Aproximação (antes da inserção) &
Dinâmica orbital/relativa ativa conforme etapa (ECI ou \gls{lvlh}); leis de controle translacionais e de atitude conforme o regime; \gls{mr} sem operação de inserção. \\
1 & Inserção &
Dinâmica acoplada espaçonave--\gls{mr} (Lagrange); cinemática inversa do \gls{mr}; cinemática direta e Jacobianos; controle PD articular do \gls{mr}; controle de atitude com compensação de reação do \gls{mr}. \\
2 & Acoplado &
Estado travado após sucesso: referência do efetuador mantida no ponto final de inserção; dinâmica e controladores mantidos conforme modo acoplado. \\
\bottomrule
\end{tabular}
\end{table}
%%%%%%%%%%%%%%%%%%

\section{Critérios de transição e sucesso no acoplamento}
\label{apD_criterios_transicao_docking}

\begin{table}[H]
\centering
\caption{Critérios numéricos de transição e de sucesso no acoplamento.}
\label{ttab_apD_criterios_docking}
\begin{tabular}{lr}
\toprule
\textbf{Parâmetro} & \textbf{Valor} \\
\midrule
Tolerância posicional $\|\vec{e}_p\|$ [m] & 0,1 \\
Tolerância angular $\theta$ [graus] & 10 \\
Comprimento máximo de inserção $s_{\max}$ [m] & 0,05 \\
Velocidade de inserção $\dot{s}$ [m/s] & 0,01 \\
Limiar de ativação por distância [m] & 10 \\
\bottomrule
\end{tabular}
\end{table}

%%%%%%%%%%%%%%%%%%%%%

\begin{table}[H]
\centering
\caption{Regras de transição entre estados do sequenciador de acoplamento.}
\label{ttab_apD_regras_transicao}
\small
\begin{tabular}{ll}
\toprule
\textbf{Transição} & \textbf{Condição lógica} \\
\midrule
0 $\rightarrow$ 1 &
$\|\vec{e}_p\| \leq 0,10$ \; e \; $\theta \leq 10$. \\
1 $\rightarrow$ 2 &
$\|\vec{e}_p\| \leq 0,10$ \; e \; $s \geq 0,05$. \\
% Reinicialização (retorno a 0) &
% Se a distância de operação excede 10 m, zera-se $s$ e retorna-se ao estado 0. \\
% \bottomrule
\end{tabular}
\end{table}

%%%%%%%%%%%%%%%%%

% =========================================================================
% APÊNDICE E — Critério de acoplamento
% =========================================================================

\section{Critério de acoplamento e detecção de sucesso}
\label{ap_apendiceE_acoplamento}

O acoplamento é declarado quando as três condições abaixo são satisfeitas simultaneamente:

\begin{itemize}
    \item $\|\vec{e}_p\| \leq 0{,}10$ m;
    \item $\theta \leq 10{,}00$ graus;
    \item $s \geq 0{,}05$ m.
\end{itemize}

A variável escalar $s$ é integrada no tempo com velocidade constante $\dot{s}=0{,}01$ m/s durante o modo de inserção.
Ao atingir $s_{\max}=0{,}05$ m, o estado passa a ser travado, mantendo a referência do efetuador no ponto final de inserção.

% \section{Condições iniciais e parâmetros orbitais por cenário}
% \label{ap_apendiceF_orbita_cenarios}

% Este apêndice consolida as variáveis orbitais utilizadas para inicialização do movimento relativo no referencial \gls{lvlh}.

% \section{Parâmetros orbitais básicos (Hill frame)}
% \label{apF_param_orbitais}

% \begin{table}[H]
% \centering
% \caption{Parâmetros orbitais utilizados no cálculo do referencial de Hill e de $\omega_{\text{orb}}$.}
% \label{ttab_apF_param_orbitais}
% \begin{tabular}{lr}
% \toprule
% \textbf{Parâmetro} & \textbf{Valor} \\
% \midrule
% $\mu_{\text{Terra}}$ [m$^3$/s$^2$] & $3,986004\times 10^{14}$ \\
% $R_{\text{Terra}}$ [m]            & $6,378137\times 10^{6}$ \\
% \midrule
% \textbf{Cenário A — altitudes} & \\
% Altitude do alvo [km]           & 600 \\
% Altitude do perseguidor [km]    & 500 \\
% \bottomrule
% \end{tabular}
% \end{table}

%%%%%%%%%%%%%%%%%%%%%%%

% \section{Condições iniciais do movimento relativo em LVLH}
% \label{apF_cond_iniciais_lvlh}

% \begin{table}[H]
% \centering
% \caption{Condições iniciais do estado relativo em \gls{lvlh}.}
% \label{ttab_apF_cond_iniciais_lvlh}
% \begin{tabular}{lr}
% \toprule
% \textbf{Variável} & \textbf{Definição} \\
% \midrule
% $\Delta \vec{R}_0$ [m] & (definida por cenário) \\
% $\Delta \vec{V}_0$ [m/s] & $[0;0;0]^T$ \\
% \bottomrule
% \end{tabular}
% \end{table}

